\section{局限性}
\label{sec:limitations}

本节把限制按"复现实验的可比性、数据、模型与评估、外部验证、因果、泛化"六层列出，并在每一条后标注本文是否已给出对应证据。凡本文未能验证者，一律标为未解决，不以推断代替证据。

\subsection{复现实验的可比性}

\begin{enumerate}
\item \textbf{三个数据集并非同分布，规模差异大。} DeepCRISPR 16\,749 条（4 细胞系、含 4 条表观通道）、Hiranniramol 1\,309 条（单细胞系、仅序列）、Labuhn 417 条（单细胞系、仅序列）。因此"信号在 Hiranniramol 上强、在 Labuhn 上无"这一差异中，\emph{规模}与\emph{数据集本身}两个因素无法完全分离。本文的辩护是：规模可以解释"更难"，但不能解释 Labuhn 上 7/7 个配置\emph{全部为负}这一方向性事实，且 Hiranniramol 与 Labuhn 在规模、细胞系数、通道数上同类却结论相反。\emph{状态：已量化（表~\ref{tab:replication}），混杂未完全消除。}
\item \textbf{每个外部数据集只有一次 \emph{single} 划分。} Hiranniramol 与 Labuhn 各只有 7 次运行（7 个模型，单次划分），因此无法给出配置间的分布或中位数区间，也不能对数据集间 $R^{2}$ 差异做显著性检验。本文的补偿措施是与官方划分无关的 5 折 CV 复核（\S\ref{sec:res-replication}）。\emph{状态：已用独立复核补偿，但重复次数仍为 1。}
\item \textbf{两个外部数据集的标签语义与 DeepCRISPR 不完全一致。} DeepCRISPR 为归一化编辑效率，Hiranniramol 原始为 \texttt{Edit Efficiency}（百分制），Labuhn 为 \texttt{KO\_reporter\_assay}。适配器统一到 $[0,1]$，但三者的测定方式、细胞体系与 readout 不同，因此跨数据集的 $R^{2}$ \textbf{不可读作"同一量的预测难度比较"}，只能读作"同一流程在不同数据集上的可预测性比较"。\emph{状态：已在方法与表注中限定。}
\item \textbf{外部数据集没有表观遗传通道，因此无法检验环境结论的复现性。} 本文关于"环境通道无增量预测证据"的结论\textbf{仅来自 DeepCRISPR}，没有第二个带表观通道的数据集可供复现。\emph{状态：未解决。}
\end{enumerate}

\subsection{数据层面}

\begin{enumerate}
\item \textbf{使用公开数据，非本研究自建实验。} 全部标签与表观轨道来自公开数据集\citep{Chuai2018GenomeBiol,Hiranniramol2020Bioinformatics,Labuhn2018NAR}，其 library 设计、实验批次与筛选条件由原始研究决定；本研究无法追溯每个位点的实验批次效应。\emph{状态：已确认的边界，不可由本文数据解决。}
\item \textbf{DeepCRISPR 内部细胞系数量有限且不平衡。} 仅 4 个细胞系，样本量 2\,076--8\,101（相差近 4 倍）。混合与 LOCO 划分下，训练与测试的细胞系构成差异会同时改变样本量与序列组成，本文无法完全分离二者影响。\emph{状态：已量化（表~\ref{tab:dataset}），未解决。}
\item \textbf{环境变量仅四条且分辨率有限。} CTCF、DNase、H3K4me3、RRBS 为逐位点二值可及性编码；四个通道不足以代表完整 cell state（三维染色质结构、转录因子占据、细胞周期、DNA 修复状态等未纳入）。部分细胞系--通道组合近似常数，会直接压缩可检测效应上限。\emph{状态：已报告效应量级与区间（\S\ref{sec:res-environment}），但未逐通道量化方差，属未解决。}
\item \textbf{泄漏审计的身份类仅覆盖序列与反向互补。} 本文审计覆盖 sequence 与 reverse-complement 两个层级，未覆盖同一 locus 的不同序列、同源但不完全相同的序列以及公共数据库层面的重复提交。\emph{状态：审计边界已明确（\S\ref{sec:res-data}）；更广层级的污染未排除。}
\item \textbf{历史泄漏批次的对照受限于两批可比的统计量。} 本文仅在 \emph{single}/\emph{mixed}/\emph{all} 三个中位 $R^{2}$ 上做前后对照（\S\ref{sec:res-leakage}），未对旧批次重算全部下游分析。\emph{状态：对照范围有限，已在文中标注。}
\end{enumerate}

\subsection{模型与评估层面}

\begin{enumerate}
\item \textbf{预测性能总体较弱，且在 Labuhn 上为负。} DeepCRISPR \emph{single} 中位 $R^{2}$ 为 $0.067$--$0.120$，\emph{mixed} 为 $0.041$--$0.094$，LOCO 为 $-0.027$ 至 $+0.010$；Labuhn 全部为负（$-0.725$ 至 $-0.050$）。本文不将预测精度作为主要贡献，也不据此判断生物学重要性。\emph{状态：已量化（表~\ref{tab:prediction}、表~\ref{tab:replication}）。}
\item \textbf{LOCO 的"分布漂移"与"训练集缩减"混杂。} 为保证无同源泄漏，LOCO 训练池剔除了与留出系共享身份类的全部序列（HCT116 留出时剔除 4\,266 条、HeLa 4\,285 条、HL60 173 条、HEK293T 0 条），训练样本量随留出系由 12\,313 降至 1\,453。跨细胞系性能下降同时包含两种原因，本文无法在本数据中定量分离。\emph{状态：剔除条数已量化，混杂未解决。}
\item \textbf{模型选择未独立于评估集。} 平台在同一批测试集上比较 7 个模型配置与 16 个环境组合，属探索性比较；虽有配对设计与 bootstrap/置换稳健性检验，仍存在选择偏差风险。本文不把"某配置最优"作为结论。\emph{状态：已声明，未通过嵌套验证解决。}
\item \textbf{不同模型的 attribution 语义不同。} 线性系数、TreeSHAP、IG、ISM 与注意力权重在定义、尺度与可加性上均不同；本文在模型内部归一化后比较相对分布，并禁止跨模型比较绝对值。\emph{状态：已在方法与图注中限定。}
\item \textbf{归因对超参数敏感。} CNN $k=7$ 的归因谱与 $k=3/5$ 差异较大（\S\ref{sec:res-kernel}），说明解释结果依赖模型配置；本文通过保留多配置交叉验证暴露该敏感性，而不选取单一"最好"配置。\emph{状态：已展示。}
\item \textbf{自助区间的推断单位是模型而非样本。} 因子级区间为 \emph{Model-Level Bootstrap Interval}（$n=7$ 个模型配置）；7 个模型共享同一份数据、标签、预处理与划分，\emph{不构成独立生物学重复}，该区间刻画 effect 的跨模型稳健性，不应解释为总体抽样置信区间，$n=7$ 下百分位自助法的覆盖性质亦无严格保证。\emph{状态：已限定。}
\item \textbf{部分 run 数值发散。} 按测试集口径，20 个 Linear Regression run 出现 $|R^{2}|\ge10$（共线性设计矩阵导致病态求解），全部被隔离而非删除，相关统计量以有效 run 数报告。\emph{状态：已量化并在表中单列。}
\end{enumerate}

\subsection{外部模型验证层面}

\begin{enumerate}
\item \textbf{外部验证只有 8 个实例。} 定向扰动的一致性（7/8）与幅度相关（$r=0.629$，$\rho=0.476$）均基于 $n=8$，本文不对其做显著性检验，也不把 $0.875$ 读作总体比例估计。\emph{状态：已限定；扩大样本量属未来工作。}
\item \textbf{CRISPRon 本身也有误差，且不是"真值"。} 本文使用 CRISPRon 作为\emph{独立系统}，而不是作为金标准。两个模型系统方向一致，只能说明它们捕捉到同一统计结构；若该结构本身是数据中的系统性偏差，两者会一起错。\emph{状态：已明确（\S\ref{sec:res-external} 解释与边界）。}
\item \textbf{WT 选择规则虽已预登记，但仍受候选池约束。} 8 条 WT 取自 DeepCRISPR \emph{single} 测试集且要求第 18 位为 C，因此该验证的总体是"第 18 位为 C 的 DeepCRISPR 序列"，不能自动外推到其它碱基或其它数据集。\emph{状态：规则与版本号已记录，适用范围已限定。}
\item \textbf{30-mer 构造依赖参考基因组与链向判定。} 外部验证需要从 hg19 补齐侧翼（本文 23\,nt 序列本身不含侧翼）。链向取窗曾出现实现错误（负链窗口反向），修正后与参考序列逐条一致；但该步骤的正确性依赖坐标与基因组版本，本文通过逐条 QC 与额外靶点计数来约束，仍属需要外部复核的环节。\emph{状态：已逐条 QC 并记录；已确认 hg19 与数据一致（hg38 为 6/6 不匹配）。}
\end{enumerate}

\subsection{因果层面}

\begin{enumerate}
\item \textbf{预测关联不等于因果关系。} 本文全部结果均为观察性数据上的预测、归因与模型反事实分析，不能支持"某因素导致编辑效率变化"的陈述。
\item \textbf{in-silico 饱和突变与 CRISPRon 打分都是模型内部反事实（E4），不是实验验证（E5）。} 第 18 位 C$\rightarrow$A 的方向性结论刻画的是\emph{模型}的敏感性，而非 Cas9 在细胞内的行为。
\item \textbf{无湿实验验证。} 本文未开展任何 guide 效率测定、竞争实验或结构实验，所有"候选"均停留在待验证状态（E5 为空）。
\item \textbf{观测性碱基关联存在混杂。} 第 18 位 C 与 A 的效率差在不同细胞系中幅度与方向不同，而四个细胞系的序列组成、样本量与标签分布均不同，"位置 18 的碱基效应"与"携带该碱基的序列子集的其它特征"在本数据中不可完全分离（E2）。
\item \textbf{模型归因的一致性不等于生物学正确性（E3）。} 多个模型类给出一致归因只能说明它们在当前数据上依赖了相似的特征，不能排除它们共同依赖了数据中的混杂结构。这是本文引入外部模型验证的原因，也是该验证仍止步于 E4 的原因。
\end{enumerate}

\subsection{泛化层面}

\begin{enumerate}
\item \textbf{结论不可直接外推到其它 Cas 蛋白。} 本文仅使用 SpCas9（NGG PAM）数据，不覆盖 Cas12a、碱基编辑器、先导编辑或其它编辑模态。
\item \textbf{结论不可直接外推到其它细胞系与实验协议。} LOCO 结果显示跨细胞系迁移失败，且三个数据集的复现结果互相矛盾，因此性能结论不应外推到未参与训练的数据集或细胞类型。
\item \textbf{"序列可预测性不可移植"这一结论本身也需在更多数据集上检验。} 本文只有两个外部复现数据集，其中一成一败。更一般的主张需要更多独立数据集的系统采样；本文只报告这一采样结果，不把它当作普遍定律。
\item \textbf{平台层面的结论比具体数值更可外推。} "流程可追溯、泄漏可审计、证据可分层、外部模型可验证"属于方法学贡献；具体效应量级与候选因素均限于本文数据集。
\end{enumerate}
